% Section propre à la version complète : détection hors FPGA
\section{Détection d'objets hors FPGA : GPU embarqués, NPU et microcontrôleurs}
\label{sec:detection-autres}

\subsection{GPU embarqués et NPU}

\paragraph{Premiers bancs d'essai : des protocoles incomplets.}
\citet{2022-feng-benchmark-analysis-yolo-performance-edge} comparent \glslink{yolov3}{YOLOv3} et \glslink{yolov4}{YOLOv4}
(complets et tiny) sur \glslink{jetson-nano}{Jetson Nano}, \glslink{xavier-nx}{Xavier NX} et Raspberry~Pi~4 + \glslink{ncs2}{NCS2}, en FP16. Leurs choix
limitent la portée des résultats :
\begin{itemize}
  \item les \glslink{jetson}{Jetson} tournent sous Darknet, sans TensorRT, choisi exprès pour rester
    comparable au Raspberry~Pi \src{§3.2} ;
  \item aucune précision n'est mesurée : la \enquote{mean confidence} est une moyenne des
    scores de confiance, sans vérité terrain \src{§4.1} ;
  \item la portée de la puissance n'est pas décrite.
\end{itemize}
Résultat marquant : après conversion par OpenVINO, \glslink{yolov3-tiny}{YOLOv3-tiny} sur \glslink{ncs2}{NCS2} tombe à 0\,\% de
confiance moyenne, contre 59,7\,\% sur les \glslink{jetson}{Jetson} \src{Table~3}. La conversion vers un
accélérateur peut dégrader la détection sans que rien ne l'indique.

\citet{arxiv-2211.11647-benchmarking-edge-computing-devices-grape} comparent un gros
détecteur, RetinaNet \glslink{resnet}{ResNet-50}, sur \glslink{tx2}{Jetson TX2}, \glslink{edgetpu}{Coral} \glslink{edgetpu}{Edge TPU}, \glslink{zcu104}{ZCU104} et \glslink{kv260}{KV260}, avec une
\glslink{map}{mAP} mesurée et une puissance à l'entrée de la carte \src{§2.5}. Les FPGA atteignent 14 à
25\,\glslink{fps}{FPS}, les GPU embarqués 3 à 5\,\glslink{fps}{FPS} \src{Abstract}. Mais la \glslink{tx2}{TX2} tourne en FP32 ou FP16
et les FPGA en INT8 : l'écart mêle l'effet du matériel et celui de la quantification. Les
\glslink{jetson-nano}{Jetson Nano} ont été exclues faute de mémoire pour la conversion \src{§4}. La plupart des
valeurs ne figurent que dans des graphiques.

\paragraph{Baselines Jetson Orin.} Deux travaux récents fournissent les baselines les plus
utiles du corpus :
\begin{itemize}
  \item \citet{arxiv-2502.15737-performance-analysis-you-only-look} mesurent \glslink{yolov8n}{YOLOv8n} et s
    en 640$\times$640 sous TensorRT (FP32, FP16, INT8), avec une puissance par télémétrie
    (tegrastats) et une énergie par inférence \src{§3.2, Table~4}. Leur jeu de données est
    privé et facile : une seule cible, en intérieur. La PTQ INT8 y coûte 6,4 à 12,3 points
    de \glslink{map5095}{mAP50-95}, alors que la discussion parle d'une perte non significative \src{§6}.
  \item \citet{arxiv-2409.16808-comprehensive-evaluation-deep-learning-object}\pp{} donnent
    la seule \glslink{map}{mAP} COCO val du corpus sur \glslink{orin-nano}{Orin Nano} et sur \glslink{hailo8}{Hailo-8}. Leur protocole est
    réutilisable : service à une requête à la fois, repos mesuré puis soustrait, trois
    répétitions \src{§III}. La précision numérique des moteurs TensorRT et Hailo n'est pas
    indiquée, et la \glslink{map}{mAP} est identique sur \glslink{jetson-nano}{Jetson Nano} et \glslink{orin-nano}{Orin Nano}, ce qui laisse douter
    qu'elle ait été mesurée sur chaque carte \src{§IV-C}.
\end{itemize}

\paragraph{Jetson et FPGA dans un même papier.} Dans \texttt{benchmarks.csv}, quatre
articles seulement ont des lignes à la fois sur FPGA et sur \glslink{jetson}{Jetson}, et aucun sur \glslink{orin}{Orin} : les
seules lignes \glslink{orin}{Orin} viennent de Rey et d'Alqahtani, sans FPGA.
\begin{itemize}
  \item \citet{arxiv-2304.11580-framework-benchmarking-real-time-embedded} mesurent le même
    \glslink{yolov4}{YOLOv4} sur \glslink{agx-xavier}{Jetson AGX Xavier} (TensorRT FP32, FP16, INT8) et sur \glslink{zcu104}{ZCU104} (\glslink{dpu}{DPU}), avec une
    \glslink{map}{mAP} COCO calculée à partir des sorties de chaque carte et une puissance de carte au
    multimètre \src{§3, §4.1}. En INT8, la \glslink{jetson}{Jetson} est plus rapide (22,0 contre 54,42\,\glslink{u-ms}{ms}
    d'inférence), plus précise (38,6 contre 37,5) et plus efficace (2,05 contre
    1,31\,\glslink{fpsw}{FPS/W}) \src{Table~1}. Réserves : un seul modèle et deux grosses cartes de
    développement ; le post-traitement de la \glslink{jetson}{Jetson} est en numpy alors que celui de la \glslink{dpu}{DPU}
    est optimisé \src{§4.2} ; des moyennes sur six passes, sans pire cas \src{§4.1}. Les puissances relatives, repos déduit, sont
    égales (7,0\,\glslink{u-w}{W}, \src{Table~1}). Le même banc recommande le mode 15\,\glslink{u-w}{W} de la \glslink{jetson}{Jetson}
    pour l'efficacité \src{§4.4}.
  \item \glslink{skynet}{SkyNet} \citep{arxiv-1909.09709-skynet-hardware-efficient-method-object}\pp{} :
    même modèle sur \glslink{tx2}{TX2} (FP32) et \glslink{ultra96}{Ultra96} (virgule fixe 9/11~bits), mesuré par les
    organisateurs de \glslink{dac-sdc}{DAC-SDC} 2019. La \glslink{tx2}{TX2} va plus vite (67,33 contre 25,05\,\glslink{fps}{FPS}) pour
    13,50 contre 7,26\,\glslink{u-w}{W}, portée non dite \src{Tables~5--6}.
  \item \citet{arxiv-2407.17647-energy-efficient-artefact-detection-accelerator}
    (\glslink{ae}{autoencodeur} de détection d'artefacts, pas un détecteur d'objets) donnent le résultat
    inverse : \glslink{zcu104}{ZCU104} (\glslink{dpu}{DPU}, INT8) en 4,0\,\glslink{u-ms}{ms} contre 10,36\,\glslink{u-ms}{ms} pour une \glslink{xavier-nx}{Xavier NX} en INT8
    \src{§IV-B, Table~VI}. Mais la \glslink{xavier-nx}{Xavier NX} est bridée à 10\,\glslink{u-w}{W} et les puissances n'ont pas
    la même portée (rails PMBus contre jtop).
  \item \citet{2024-aggarwal-chameleon-dual-memory-replay-online} comparent un
    apprentissage continu sur \glslink{zcu102}{ZCU102} (413\,\glslink{u-ms}{ms} par image d'entraînement) et sur \glslink{jetson-nano}{Jetson Nano}
    (33\,\glslink{u-ms}{ms}), mais à batch~1 contre batch~8 \src{Table~VI} (section~\ref{sec:cl-fpga}).
\end{itemize}
Les baselines GPU des autres travaux de détection sur FPGA sont faibles ou reprises :
\glslink{jetson-nano}{Jetson Nano} en FP32 plutôt qu'en INT8
\citep[§5.3]{2024-zhang-end-end-acceleration-yolo-object}, CPU et GPU de serveur repris
d'un autre article \citep[Table~3]{2024-yan-fpga-yolov5-accelerator-real-time}, GPU de
bureau \citep[Table~10]{2026-fata-low-cost-fpga-enhanced-cnn}. Dans les deux articles où une
\glslink{jetson}{Jetson} exécute le même détecteur que le FPGA (Schlosser et al.\ et \glslink{skynet}{SkyNet}), elle est plus
rapide ; elle est aussi plus efficace en \glslink{fpsw}{FPS/W} (puissance de carte chez Schlosser, de
portée non dite chez \glslink{skynet}{SkyNet}) \deduit.

\subsection{Microcontrôleurs}

Sur MCU, la contrainte principale est la mémoire. Le survey de
\citet{2025-el-zeinaty-designing-object-detection-models-tinyml} recense des détecteurs de
55 à 789\,MMACs, avec des pics de SRAM de 53 à 511\,\glslink{u-ko}{kB} et des \glslink{map}{mAP} VOC de 51 à 75\,\%
\src{Table~8}, chiffres repris des articles d'origine. Ce tableau ne donne ni latence ni
énergie, bien que le survey présente ces métriques comme essentielles \src{§4.2}.

\citet{arxiv-2110.15352-mcunetv2-memory-efficient-patch-inference} exécutent les premiers
blocs du réseau par patchs, ce qui réduit le pic de SRAM, et cherchent conjointement
l'architecture et l'ordonnancement \src{§3}. MCUNetV2-M4 atteint 64,6\,\% de \glslink{map}{mAP} VOC en
247\,\glslink{u-ko}{kB} de SRAM sur un \glslink{cortex-m4}{Cortex-M4} \src{Table~3}. Ni la latence ni l'énergie de ces
détecteurs ne sont données.

\citet{arxiv-2306.00001-tinyissimoyolo-quantized-low-memory-footprint} mesurent un petit
\glslink{yolo}{YOLO} en 88$\times$88 sur quatre MCU, avec la puissance du MCU seul \src{§III}. Sur le
MAX78000, qui intègre un accélérateur CNN, une inférence dure 5,5\,\glslink{u-ms}{ms} et coûte 196\,\glslink{u-uj}{µJ} ; le
\glslink{cortex-m7}{Cortex-M7} d'un \glslink{stm32h7a3}{STM32H7A3} met 359\,\glslink{u-ms}{ms} \src{§III, Fig.~2}$^{m}$. L'accélérateur intégré fait
donc la différence (plus de 65 fois plus rapide). La \glslink{map}{mAP} de 58,5\,\% porte sur 3 classes
de VOC et au plus 3 objets par image \src{Table~II} ; elle ne se compare pas aux \glslink{map}{mAP}
de \glslink{mcunetv2}{MCUNetV2} sur 20 classes. Le papier lui-même donne trois valeurs différentes selon
l'endroit (voir la fiche).

Aucun article du corpus ne mesure un même détecteur sur MCU et sur FPGA. Le détecteur FPGA
le plus proche d'un MCU en taille est le \glslink{lp-yolo}{LP-YOLO Slim} de
\citet{arxiv-2609.24757-npu-accelerator-quantized-real-time}\pp{} : environ 0,35\,M de
paramètres pour \glslink{lp-yolo}{LP-YOLO} \src{§III-A}, poids sur 2~bits, 2,91\,\glslink{u-w}{W} de carte sur \glslink{zynq}{Zynq}
\glslink{xc7z020}{XC7Z020} \src{§V-D}. Il n'a pas été porté sur MCU, et ses 2,91\,\glslink{u-w}{W} sont une puissance de
carte, alors que \glslink{tinyissimoyolo}{TinyissimoYOLO} mesure le MCU seul : un portage des deux côtés, à portée
d'énergie égale, manque.

\subsection{Tableau des détecteurs hors FPGA}

\begin{table}[!htbp]
\centering\footnotesize
\caption{Détecteurs hors FPGA dans \texttt{benchmarks.csv} (sélection). Précisions non
comparables d'une ligne à l'autre. $^{m}$ : confiance \emph{medium}. Portée : c =
\emph{chip}, b = \emph{board}, u = \emph{unspecified}.}
\label{tab:detection-autres}
\begin{tabularx}{\textwidth}{@{}P{1.8cm}P{2.6cm}LP{2.5cm}rrrl@{}}
\toprule
Travail & Modèle, précision & Plateforme & Métrique & Préc. & \glslink{fps}{FPS} & ms & W \\
\midrule
Rey & \glslink{yolov8n}{YOLOv8n} 640 INT8 & \glslink{orin-nano}{Orin Nano} & \glslink{map5095}{mAP50-95}, jeu privé & 71,9 & 44,9 & 23,16 & 7,48 (c) \\
Rey & \glslink{yolov8n}{YOLOv8n} 640 FP16 & \glslink{orin-nano}{Orin Nano} & \glslink{map5095}{mAP50-95}, jeu privé & 84,4 & 40,0 & 26,70 & 7,42 (c) \\
Rey & \glslink{yolov8n}{YOLOv8n} 640 INT8 & \glslink{orin-nx}{Orin NX} & \glslink{map5095}{mAP50-95}, jeu privé & 71,9 & 65,83 & 15,16 & 10,31 (c) \\
Rey & \glslink{yolov8n}{YOLOv8n} 640 FP32 & RPi~5 (CPU) & \glslink{map5095}{mAP50-95}, jeu privé & 84,4 & 8,47 & 118 & 5,44 (c)$^{m}$ \\
Alqahtani & \glslink{yolov8n}{YOLOv8n} 640, précision \nr & \glslink{orin-nano}{Orin Nano} & \glslink{map}{mAP} COCO val & 31 & \nr & 16 & \nr$^{m}$ \\
Alqahtani & \glslink{yolov8n}{YOLOv8n} 640, précision \nr & RPi~5 + \glslink{hailo8}{Hailo-8} & \glslink{map}{mAP} COCO val & 32 & \nr & 76 & \nr$^{m}$ \\
Schlosser & \glslink{yolov4}{YOLOv4} 512 INT8 & \glslink{agx-xavier}{AGX Xavier} & \glslink{map5095}{mAP50-95} COCO & 38,6 & \nr & 22,0 & 16,5 (b) \\
Schlosser & \glslink{yolov4}{YOLOv4} 512 FP16 & \glslink{agx-xavier}{AGX Xavier} & \glslink{map5095}{mAP50-95} COCO & 39,8 & \nr & 35,32 & 20 (b) \\
\glslink{skynet}{SkyNet} & FP32 & \glslink{tx2}{TX2} & \glslink{iou}{IoU}, 1 objet & 73,1 & 67,33 & \nr & 13,5 (u) \\
Feng & \glslink{yolov3-tiny}{YOLOv3-tiny} 416 FP16 & \glslink{xavier-nx}{Xavier NX} (sans TensorRT) & confiance moy. & 59,7 & 41,1 & \nr & 13,5 (u) \\
Feng & \glslink{yolov3-tiny}{YOLOv3-tiny} 416 FP16 & RPi~4 + \glslink{ncs2}{NCS2} & confiance moy. & 0 & 18,8 & \nr & 6,5 (u)$^{m}$ \\
\glslink{tinyissimoyolo}{Tinyissimo\-YOLO} & 88$\times$88 INT8 & MAX78000 & \glslink{map}{mAP} VOC 3 classes & 58,5 & 180 & 5,5 & \nr$^{m}$ \\
\glslink{tinyissimoyolo}{Tinyissimo\-YOLO} & 88$\times$88 INT8 & \glslink{stm32h7a3}{STM32H7A3} & \glslink{map}{mAP} VOC 3 classes & 58,5 & \nr & 359 & \nr$^{m}$ \\
\glslink{mcunetv2}{MCUNetV2} & M4 + \glslink{yolov3}{YOLOv3} INT8 & \glslink{stm32f412}{STM32F412} & \glslink{map}{mAP} VOC & 64,6 & \nr & \nr & \nr \\
\glslink{mcunetv2}{MCUNetV2} & H7 + \glslink{yolov3}{YOLOv3} INT8 & \glslink{stm32h743}{STM32H743} & \glslink{map}{mAP} VOC & 68,3 & \nr & \nr & \nr \\
\bottomrule
\end{tabularx}
\end{table}

Sources : Rey \citep[Table~4, §4]{arxiv-2502.15737-performance-analysis-you-only-look} ;
Alqahtani \citep[§IV-B--C]{arxiv-2409.16808-comprehensive-evaluation-deep-learning-object} ;
Schlosser \citep[Table~1]{arxiv-2304.11580-framework-benchmarking-real-time-embedded} ;
\glslink{skynet}{SkyNet} \citep[Table~5]{arxiv-1909.09709-skynet-hardware-efficient-method-object}\pp{} ;
Feng \citep[Table~3]{2022-feng-benchmark-analysis-yolo-performance-edge} ; \glslink{tinyissimoyolo}{TinyissimoYOLO}
\citep[§III, Table~II]{arxiv-2306.00001-tinyissimoyolo-quantized-low-memory-footprint} ;
\glslink{mcunetv2}{MCUNetV2} \citep[Table~3]{arxiv-2110.15352-mcunetv2-memory-efficient-patch-inference}. Chez
Rey et al., la portée \emph{chip} correspond à une télémétrie embarquée (rails du module
\glslink{jetson}{Jetson}, PMIC du Raspberry~Pi), pas à un wattmètre sur la carte. La latence de Schlosser est
l'inférence seule ; la mesure de \glslink{skynet}{SkyNet} est faite par les organisateurs de \glslink{dac-sdc}{DAC-SDC}.
